\chapter{Conclusion and future work}
\label{chap:conclusion}

\section{Summary of contributions}
\label{sec:conc-summary}

This project tackled the problem of detecting and classifying
radiological anomalies in noisy gamma-dose time series. Starting from
four months of real 2015 recordings, we built a complete pipeline from
data to deployment-ready models:

\begin{enumerate}[leftmargin=*, itemsep=0.25em]
    \item A \textbf{noise model} was fitted on the real signal --- a
    high-frequency Gaussian component on top of an Ornstein--Uhlenbeck
    baseline drift --- and its parameters were extracted automatically
    by a calibration script.
    \item Six \textbf{anomaly templates} were designed jointly with the
    supervisor as a structured CSV format (keypoints, allowed
    deformation quadrants, interpolation modes), making it
    straightforward to add new shapes in the future.
    \item A scalable \textbf{synthetic dataset generator} produces a
    ten-million-point labelled signal in constant memory, mixing the
    noise model with random deformations of the six templates.
    \item Three \textbf{neural network architectures} were implemented
    in PyTorch with a shared interface, ranging from a compact
    $23\,000$-parameter 1D-CNN to a $110\,000$-parameter hybrid
    CNN+Transformer model with multi-head self-attention.
    \item A \textbf{sliding-window inference engine} and an
    \textbf{event-level evaluation protocol} bridge the gap between
    sample-level classification accuracy and the operational metric
    that matters (do we detect each anomaly event, and do we get its
    class right?).
    \item A \textbf{robustness study} sweeps background noise and shape
    deformation independently and reveals the gracious-degradation
    properties of each architecture.
    \item A \textbf{real-data deployment trial} on the original 2015
    recordings produced an honest negative result --- the models do
    not work well out of the box --- followed by a careful diagnostic
    that identifies five distinct causes of the failure and a punch-list
    of seven concrete remediations.
\end{enumerate}

\section{Architecture recommendation}
\label{sec:conc-archi}

On the synthetic benchmark the ResNet 1D is the clear winner: it
achieves a perfect $100\%$ event-level macro F1 on the test segment,
degrades the most gracefully under both noise and deformation stress,
and remains compact enough ($\sim 75\,000$ parameters) to retrain in
under ten minutes on a laptop GPU. The CNN+Attention model is
competitive on synthetic data but exhibits a more severe
out-of-distribution collapse on real data; we keep it in the pipeline
mainly because it offers the most promising path to longer-window
analyses (\cref{sec:future-attention}). The baseline 1D-CNN is the
fastest to train but its catastrophic drop under high noise makes it
unfit for a beacon that must run across seasons.

If a single architecture must be picked for the next iteration, the
recommendation is \textbf{ResNet 1D}, possibly with a binary
detection stage in front of it (\cref{sec:real-remediation}, R6).

\section{Limitations}
\label{sec:conc-limits}

Several known limitations of the current pipeline are worth recording
explicitly.

\begin{itemize}[leftmargin=*, itemsep=0.25em]
    \item \textbf{Distribution shift.} As detailed in
    \cref{chap:realdata}, the synthetic training distribution does not
    match the deployment distribution along several dimensions (class
    prior, anomaly duration, baseline drift, intermediate signals).
    The system as it stands cannot be deployed on raw beacon data
    without first addressing these mismatches.
    \item \textbf{No real ground truth.} The four months of real data
    have no expert labelling, so we have no quantitative measure of
    how the models do on real anomalies. The qualitative inspection
    in \cref{sec:real-symptoms} is the best we can do for now.
    \item \textbf{Single beacon.} The pipeline assumes a univariate
    time series. Real surveillance networks involve several beacons
    and several covariates (temperature, humidity, pressure).
    Exploiting these as additional input channels is left for
    future work.
    \item \textbf{Black-box decisions.} The trained models output a
    class label but not a justification. The project brief mentions
    Grad-CAM \cite{selvaraju2017gradcam} as a candidate visualisation
    method; the architectures expose the relevant hooks
    (\texttt{last\_conv\_features}) but the Grad-CAM analysis itself
    was outside the time budget of this iteration.
\end{itemize}

\section{Future work}
\label{sec:conc-future}

The seven remediation steps of \cref{sec:real-remediation} are the
priority list for the next iteration. Beyond them, three longer-term
directions are worth flagging.

\subsection{Self-supervised pre-training on the real signal}

A natural way to bridge the synthetic--real gap is to pre-train the
convolutional stem on the four (and ideally more) months of real,
unlabelled data, using a self-supervised objective such as
masked-segment reconstruction. The stem then learns features adapted
to the real signal statistics; only the classifier head needs to be
fine-tuned on the synthetic dataset. This approach has been very
successful in audio and biomedical time-series and is directly
applicable here.

\subsection{Longer windows and stronger attention}
\label{sec:future-attention}

The CNN+Attention model did not benefit from its long-range mixing on
$512$-sample windows because the templates rarely extend past
$\sim 800$ samples. With longer windows ($W = 4096$ or more) the
attention model could in principle disambiguate close events
(\texttt{M} versus two consecutive \texttt{bell}s), reason about
context (a slow ramp during which a sharp spike occurs), and detect
multi-event episodes. Scaling the attention up requires the usual care
(longer sinusoidal encoding, more heads or layers, possibly
linear-complexity attention variants) but the architecture is already
in place.

\subsection{Explainability with Grad-CAM and attention visualisation}

The two convolutional architectures expose their last convolutional
layer as \texttt{last\_conv\_features}, ready to be hooked by a
Grad-CAM implementation. For the attention model, the attention
matrices themselves are interpretable saliency maps over the token
sequence. Adding both to the inference scripts would let the supervisor
visually verify that the model is paying attention to the right part
of an event when it predicts a class, which is the kind of
explainability the project brief explicitly requested.

\subsection{Multi-beacon and multi-modal extensions}

A practical surveillance system never works with a single signal in
isolation; pressure, temperature, humidity and the gamma counts of
neighbouring beacons are all informative. Extending the input to a
small number of channels is mechanical (the convolutions are already
multi-channel by design); the harder problem is collecting a labelled
multi-modal dataset, which the synthetic generator could be extended
to produce.

\section{Final words}
\label{sec:conc-final}

The project ends in an honest mixed state. On the benchmark we
designed for ourselves the system works very well --- all three
architectures cross the $90\%$ event-level F1 line and the residual
network is essentially perfect. On the real data it does not, and the
gap between the two is dominated by a small number of identifiable
mismatches between the synthetic training distribution and the
deployment distribution. Each of those mismatches has a known
remediation; the next iteration of the project will apply them and
re-run the same evaluation pipeline that is already in place. Many of
the pieces required for the next round --- the data generator, the
event-level evaluator, the robustness study, the three architectures
--- are reusable as-is.

The single most useful lesson of this iteration is therefore not any
particular F1 number but a methodological one: \emph{do not benchmark
a deep model on the same kind of data it was trained on, and call it
done}. The first contact with real data is what revealed where the
pipeline really stands, and the diagnostic infrastructure we built to
make that contact diagnosable is, in the long run, the most valuable
output of the project.
